\scanpage{37}
La construction du \S~I nous fournit donc des éléments remarquables
\[
\left\{\begin{aligned}
r_0,r_1,r_\infty&\in\mathrm{Gl}(2,\mathbb Z)\\
\lambda_0,\lambda_1,\lambda_\infty&\in\mathrm{Sl}(\unsure{1|2},\unsure{2|\mathbb Z})
\end{aligned}\right.
\]
satisfaisant les relations remarquables
\[
\left\{\begin{aligned}
&r_0^2=r_1^2=r_\infty^2=1\\
&\lambda_0=r_\infty r_1,\ \lambda_1=r_0r_\infty,\ \lambda_\infty=r_1r_0
&&\text{d'où }\lambda_\infty\lambda_1\lambda_0=1\\
&\lambda_\infty^{\deleted{6}3}=\lambda_1^2=\omega\ \text{(matrice $-1$)},&&\text{d'où }\lambda_\infty^6=\lambda_1^{\unsure{4|1}}=1\,;
\end{aligned}\right.
\]
sans impliquer pourtant que les hom correspondants
\[
\left\{\begin{aligned}
G^\circ_{\deleted{\unsure{3}}\infty,3}&\longrightarrow\mathrm{Sl}(2,\mathbb Z)'\\
G_{\infty,3}&\longrightarrow\mathrm{Gl}(2,\mathbb Z)'
\end{aligned}\right.
\]
sont des isomorphismes, i.e. sans impliquer
\[
\left\{\begin{aligned}
\mathrm{Sl}(2,\mathbb Z)'&\simeq\{\dot{\lambda}_1,\dot{\lambda}_\infty\mid\dot{\lambda}_1^2=\dot{\lambda}_\infty^3=1\}\\
\mathrm{Gl}(2,\mathbb Z)'&\simeq\{\dot{r}_0,\dot{r}_1,\dot{r}_\infty\mid
\dot{r}_0^2=\dot{r}_1^{\unsure{7|2}}=\dot{r}_\infty^2=(r_0r_\infty)^2=(r_1r_0)^3=1\}
\end{aligned}\right.
\]
Le fait qu'il en est bien ainsi \deleted{\unsure{rés}}\unsure{est une} identité remarquable, dont je ne connais pas de démonstration par voie purement algébrique. Je peux l'énoncer en disant que la structure 2-polyédrale triangulée orientée régulière, \deleted{déf} $SP_{\deleted{\unsure{\Lambda}}M\unsure{=}}$, définie par un $\mathbb Z$-module \deleted{\unsure{de}} libre de rg 2 muni d'une base symplectique, est \uline{universelle}, en ce sens
% archive sheet number: 36
